\documentclass[11pt]{article}
\usepackage[margin=27mm]{geometry}
\usepackage{amsmath,amssymb}
\usepackage[hidelinks]{hyperref}
\title{Pleniluner's Space Station\\Rendering Notes}
\author{Pleniluner}
\date{13 November 2026}
\begin{document}
\maketitle

\section{A stationary illustration}
The hero illustration has no scroll-dependent transformation. Its local image
transform is the identity, and its clipping region is the rectangular viewport:
\[
 T_{\mathrm{image}}(s)=I,\qquad C(s)=[0,w]\times[0,h]
 \quad\hbox{for every scroll position }s.
\]
Cover sizing still follows the aspect ratio of the hero container. The separate
star layer does not change the image, its bounds, or the layout of any text.

\section{Bounded drawing resolution}
For a canvas with CSS width $w$, height $h$, and device pixel ratio $d$, the
JavaScript driver selects
\[
 r=\min\left(d,1.5,\frac{1600}{\max(w,h)}\right),\qquad
 (W,H)=\left(\max(1,\operatorname{round}(rw)),
              \max(1,\operatorname{round}(rh))\right).
\]
Thus the longest buffer edge is at most 1600 pixels, apart from integer
rounding. Fragment coordinates are converted back to CSS coordinates using
$\mathbf{p}=\mathbf{x}_{\mathrm{frag}}/r$. Star spacing therefore remains
consistent across ordinary and high-density displays. No high-resolution image
is sampled by the shader.

\section{Sparse stars and smooth edges}
The plane is divided into cells of side length $64$ CSS pixels. Let
\[
 \mathbf{c}=\lfloor\mathbf{p}/64\rfloor,\qquad
 \mathbf{q}=64\,\operatorname{fract}(\mathbf{p}/64),\qquad
 H(\mathbf{c})=\operatorname{fract}\left(
 \sin(127.1c_x+311.7c_y)\,43758.5453\right).
\]
This deterministic visual hash is not a cryptographic random function.
The center and radius of a candidate star are
\begin{align*}
 \boldsymbol{\mu}&=64\left[0.175+0.65
   \begin{pmatrix}H(\mathbf{c})\\H(\mathbf{c}+7)\end{pmatrix}\right],\\
 \rho&=0.45+0.60H(\mathbf{c}+3).
\end{align*}
Vector additions use the same scalar in both coordinates, matching GLSL.
Only cells satisfying $H(\mathbf{c}+11)\geq0.84$ are visible. With
$S(a,b,x)=\operatorname{smoothstep}(a,b,x)$, their opacity is
\[
 \alpha=0.32\,\mathbf{1}_{H(\mathbf{c}+11)\geq0.84}
 \left[1-S\left(0.25\rho,\rho+0.8/r,
                  \lVert\mathbf{q}-\boldsymbol{\mu}\rVert\right)\right].
\]
Premultiplied output is $(\alpha\mathbf{k},\alpha)$ with
$\mathbf{k}=(0.72,0.78,0.75)$. This keeps transparent pixels compatible with
browser composition over the original illustration.

\section{Drawing and recovery}
The layer draws on initialization, container or window resize, and restored
WebGL context. Requests within a frame are coalesced. There is no time uniform
and no continuous animation loop. When the context is lost, GPU object
references are discarded; on restoration the program and buffer are rebuilt.
Component disposal removes observers and listeners and deletes live resources.
If context creation or shader compilation fails, CSS supplies a static texture.
The reduced-motion preference never introduces an alternative animated path.

\section{Other illustrations}
Full-frame illustrations use their intrinsic aspect ratios without scroll
scaling. Remaining background scenes retain the existing smooth arrival
function $f(t)=t^2(3-2t)$ for $0\leq t\leq1$. These scene parameters have no
effect on the hero image or the static star layer.

\section{Source map}
\begin{itemize}
 \item TypeScript: component mounting, page structure and interaction state.
 \item JavaScript: WebGL resource lifecycle and build tools.
 \item GLSL: vertex placement and per-pixel star opacity.
 \item CSS: composition, responsive layout and fallback texture.
 \item TeX: the equations and rendering constraints in this document.
\end{itemize}
\end{document}
